\documentclass[11pt]{article}
\usepackage{fontspec}
\setmainfont{Times New Roman}
\setsansfont{Arial}
\setmonofont{Consolas}
\usepackage{amsmath,amssymb,mathtools}
\usepackage{geometry}
\usepackage{microtype}
\geometry{a4paper,margin=2.05cm}
\setlength{\parindent}{1.1em}
\setlength{\parskip}{0pt}
\newcommand{\foreign}[1]{#1}
\providecommand{\Norm}{\operatorname{N}}
\providecommand{\Sp}{\operatorname{Sp}}
\providecommand{\Gg}{\mathfrak G}
\providecommand{\Hh}{\mathfrak H}
\providecommand{\Ii}{\mathfrak I}
\providecommand{\PP}{\mathfrak P}
\providecommand{\oo}{\mathfrak o}
\providecommand{\glqq}{``}
\providecommand{\grqq}{''}
\emergencystretch=220em
\allowdisplaybreaks
\sloppy
\hbadness=10000
\vbadness=10000
\hfuzz=5pt
\vfuzz=12pt
\raggedbottom
\begin{document}

\section*{42. Разложене скрежене продукты и јих максималне порјадкы}

\begin{center}
\emph{\foreign{Actualités scientifiques et industrielles 148 (1934), S. 5--15}}
\end{center}

В најпростејшем случају разложених скрежених продуктов,\footnote{Теорија скрежених продуктов, в связи с једноју мојеју лекцијеју, јест изложена в одделу II у Х. Хассе, \foreign{\emph{Theory of cyclic algebras over an algebraic number field}}, Транс. Амер. Math. Soc. 34 (1932). Тут употребјам само прве основне појмы. Брандтове основы у Хассе и в следујучем оквиру былы наново обосноване.} то јест тых, кторе дају разложене алгебры и зато имајут факторне системы асоциоване с јединицоју, можно односеньа в великој мери експлицитно следити. То остаје правдиво и при не-Галоисовых разпадных польах, то јест при максималных комутативных подпольах. Сем належече просте алгебраичне факты развијајут се в §1. В §2, при алгебраичном числовом польу како основном польу, давам експлицитне представјеньа всех максималных порјадков и јих идеалов посредством модулов и комплементарных модулов основного разпадного польа \(k\).\footnote{Та представјеньа была ми давно знана; побудку к следујучим разсмотреньам дала ми објава једного теорема Х. Хассе.}

Далеј давам разделбу на области: в једну област рачут се все таке максималне порјадкы, кторе имајут једнаковы пресек с \(k\). Тым области једно-једнозначно одповедајут порјадкам из \(k\); главному порјадку јест приредјена главна област. Максималне порјадкы једнеј области трансформујут се једне в друге идеалами, образоваными из модулов одповедајучего порјадка; особливо в главној области иде о разширеньа идеалов из \(k\). В §3 преходим к произвольным простым алгебрам. Уточньеньем разделбы на области, тако же поред пресека морајут совпадати и \(p\)-адичне компоненты максималных порјадков на разветвјених местах, можно пренести теоремы добите в разложеном случају. Особливо и тогда максималне порјадкы главној области преходет једне в друге трансформацијеју с разширеньами идеалов из \(k\).

Теоремы о главној области находет се при Цхеваллеы и Хассе;\footnote{C. Цхеваллеы, \foreign{Sur certains idéaux d'une algèbre simple}, Abh. Math. Сем. Хамбург, 1934; Х. Хассе, \foreign{Über gewisse Ideale in einer einfachen Algebra}, работа, ктора предходи овој.} Цхеваллеы такодје показује, же без уточньене разделбы, односне к разветвјеным местам, теоремы в обчем случају веч не држе.

\subsection*{§I. Матричне јединице разложених скрежених продуктов и комплементарне базы разпадных поль}

В тутом параграфу \(\Omega\) значи произвольно основно полье. Нехај \(k/\Omega\) јест Галоисово сепарабилно разширенье \(n\)-того степена, \(\Gg\) Галоисова група с елементами \(S,T,\ldots\), а \(u_S,u_T,\ldots\) одповедајучи операторы. Разгледајут се скрежене продукты с факторноју системоју јединица, теды разложене алгебры
\[
        K=\Gg\times k=u_{S_1}k+\cdots+u_{S_n}k,\qquad
        u_Su_T=u_{ST},
\]
с
\[
        z u_S=u_S z^S \qquad(z\in k).
\]
Ако положимо
\[
        E_\Gg=\sum_{S\in\Gg}u_S,
\]
из факторној системы јединица добивајут се односеньа
\[
        E_\Gg u_S=u_S E_\Gg=E_\Gg,                            \tag{1}
\]
\[
        zE_\Gg=\sum_Su_Sz^S,                                  \tag{2a}
\]
и
\[
        E_\Gg z E_\Gg
        =E_\Gg\sum_Su_Sz^S
        =E_\Gg\,\Sp(z),                                      \tag{2b}
\]
где \(\Sp(z)\) значи след од \(z\) како елемента \(k/\Omega\).

\paragraph{Теорем.}
Ако \(a_1,\ldots,a_n\) јест база \(k/\Omega\), а \(\bar a_1,\ldots,\bar a_n\) к тому комплементарна база, тогда
\[
        c_{ik}=\bar a_iE_\Gg a_k
\]
дава систем матричных јединиц в \(K\). \(K\) можно представити како модулны продукт в форме
\[
        K=kE_\Gg k
          =k\Bigl(\sum_Su_S\Bigr)k
          =\sum_{i,k}(\bar a_iE_\Gg a_k)\,\Omega .
\]
Бо по (2б) наставаје
\[
        c_{ij}c_{hk}
        =\bar a_iE_\Gg a_j\bar a_hE_\Gg a_k
        =\bar a_iE_\Gg a_k\,\Sp(a_j\bar a_h)
        =c_{ik}\,\Sp(a_j\bar a_h).
\]
То јест равно својство матричных јединиц, заради дефиницијных равнань
\[
        \Sp(a_j\bar a_h)=\delta_{jh}
\]
комплементарных баз.

\paragraph{Заметка 1.}
Ако \(Z\) јест како-коли комутативно разширјуче полье над \(\Omega\), все остаје, когда \(a_i\) значи базу \(k_Z/Z\), а \(K_Z\) ставаје се директноју сумоју поль. Субституције групы тогда сут дефиниране тым, же на \(Z\) индуцирајут идентитет.

\paragraph{Заметка 2.}
То, же \(K=kE_\Gg k\), следује такодје непосредно из того, же права страна јест ненуловы идеал простого система \(K\).

Нехај ныне \(\Omega\) има карактеристику нула. Представјенье \(K=kE_\Gg k\) разложене алгебры остаје правдиво и тогда, когда \(k\) јест не-Галоисово разпадно полье за \(K/\Omega\). Преходи се к Галоисовому разпадному польу \(\bar k\). Нехај \(\bar k\) буде Галоисово полье, одповедајуче к \(k\), \(\bar\Gg\) јего група, \(\bar u_S\) одповедајучи операторы, а \(\bar K=\bar kE_{\bar\Gg}\bar k\) разложены скрежены продукт. Подполье \(k\) одповедаје подгрупе \(\Hh\) порјадка \(h\). Положимо
\[
        \bar\Gg=\Hh S_1+\cdots+\Hh S_n
              =S_1^{-1}\Hh+\cdots+S_n^{-1}\Hh,
        \qquad
        E_\Hh=\frac1h\sum_{H\in\Hh}u_H,
\]
и
\[
        E_\Gg=\frac1h E_{\bar\Gg}
              =\sum_i u_{S_i},\qquad
        u_{S_i}=E_\Hh\bar u_{S_i}.
\]
Тогда особливо идемпотент \(E_\Hh\) твори јединично представјенье \(\Hh\), а \(E_\Gg\), тако како \(E_{\bar\Gg}\), твори јединично представјенье \(\bar\Gg\). За тако дефиниране \(E_\Gg\) и \(u_{S_i}\) односенье (1) држи по једнеј стране, а (2) држи полно. Најпрво
\[
        E_\Gg=E_\Gg E_\Hh=E_\Hh E_\Gg=E_\Hh E_\Gg E_\Hh,       \tag{3}
\]
\[
        zE_\Hh=E_\Hh z,                                      \tag{4}
\]
и следовательно
\[
        z u_{S_i}=u_{S_i}z^{S_i}.                            \tag{5}
\]
Из (5) следује (2а) сумованьем, из того (2б), а (1) следује по једнеј стране:
\[
        E_\Gg u_{S_i}=E_\Gg .
\]
Доказ \(K=kE_\Gg k\) иде ныне точно тако како выше. Како полно матрично колцо над \(\Omega\), \(K\) всебује всако полье \(n\)-того степена како подполье, особливо и \(k\).

S погледом на §3 треба јеште заметити, же преход од \(K\) к \(\bar K\) јест можны и когда \(K\) не разклада се. Поред \(k\) тогда јест и \(\bar k\) разпадно полье, теды \(K\) јест подобна к \(\bar K\), скреженому продукту из \(\bar\Gg\) с \(\bar k\). При том, бо \(k\) јест разпадно полье, подгрупе \(\Hh\), одповедајучеј к \(k\), одповедаје факторна система јединица; идемпотент \(E_\Hh\) теды јест дефиниран како выше. Из того обратно следује, же \(K\) јест изоморфна к \(E_\Hh\bar K E_\Hh\). Та алгебра јест изоморфна к колцу автоморфизмов левого идеала \(\bar K E_\Hh\), теды подобна к \(\bar K\); ранг совпада с рангом \(K\). В разложеном случају добивамо
\[
        E_\Hh\bar K E_\Hh
        =E_\Hh\bar kE_{\bar\Gg}\bar kE_\Hh
        =(E_\Hh\bar kE_\Hh)E_{\bar\Gg}(E_\Hh\bar kE_\Hh)
        =kE_\Gg k .
\]

\subsection*{§II. Максималне порјадкы разложених скрежених продуктов и јих разделба на области}

Од ныне нехај \(\Omega\) јест конечно алгебраично числово полье, \(K\) разложена алгебра \(n\)-того степена над \(\Omega\), а \(k\) Галоисово или не-Галоисово максимално комутативно подполье, теды \(K=kE_\Gg k\). Нехај \(a,\bar a\) значи пару комплементарных конечных \(\oo\)-модулов из \(k\), при чем \(\oo\) јест главны порјадок од \(\Omega\), а \(o\) главны порјадок од \(k\). Структура максималных порјадков из \(K\) релативно к \(k\) описана јест следујучими теоремами.

\paragraph{Теорем 1.}
Модулны продукт
\[
        \mathcal O_a=\bar a E_\Gg a
\]
јест максималны порјадок из \(K\); особливо
\[
        \mathcal O=\mathcal O_o=oE_\Gg o
\]
ставаје се максималным порјадком.

\paragraph{Теорем 2.}
Трансформација од \(\mathcal O\) в \(\mathcal O_a\) деје се посредством левого идеала
\[
        \mathcal L=oE_\Gg a
\]
и к ньему реципрочного правого идеала
\[
        \mathcal R=\bar a E_\Gg o
\]
по \(\mathcal O\).

\paragraph{Заметка.}
Совпаданье двух комплементарных модулов \(a,\bar a\) в продукту \(\mathcal L\mathcal R\) веде равно заради следотворјеньа, учиньенохо чрез \(E_\Gg\), к реципрочној связи.

\paragraph{Теорем 3.}
Все леве и праве идеалы по \(\mathcal O\) сут изчерпане тыми, кторе указане сут в теорему 2. Максималне порјадкы \(\mathcal O_a\) теды изчерпају всу совокупност максималных порјадков из \(K\). Кака-коли два максималне порјадкы \(\mathcal O_a\) и \(\mathcal O_b\) преходет једне в друге трансформацијеју с
\[
        \mathcal L_{ab}=\bar aE_\Gg b,
        \qquad
        \mathcal R_{ba}=\bar bE_\Gg a
\]
и те \(\mathcal L_{ab}\) и \(\mathcal R_{ba}\) изчерпају леве и праве идеалы по \(\mathcal O_a\).

\paragraph{Теорем 4.}
Когда \(\mathcal O_a\) бежи чрез все максималне порјадкы из \(K\), пресек
\[
        [\mathcal O_a,k]
\]
бежи чрез все порјадкы из \(k\); тут пресек јест всакы раз равны порјадку од \(a\). Ако все максималне порјадкы, належече к једному и тому самому порјадку в \(k\), собирајут се в једну област, тогда трансформација внутри области деје се левыми идеалами \(\bar aE_\Gg b\) по \(\mathcal O_a\) и реципрочными правыми идеалами, при чем \(a\) и \(b\) сут модулы одповедајучего порјадка в \(k\).

\paragraph{Теорем 4а.}
Особливо трансформација максималных порјадков главној области, то јест области всех максималных порјадков, всебујучи главны порјадок \(o\) од \(k\), деје се идеалами
\[
        aE_\Gg b,
\]
где \(a\) и \(b\) сут идеалы из \(k\). Другыми словами: ако \(\mathcal L\) јест левы идеал максималного порјадка главној области и \(o\) јест всебовано в правом порјадку од \(\mathcal L\), тогда \(\mathcal L\) јест разширенье \(o\)-идеала.

Доказы опирајут се на преход к појединым местам. Достаточно јест всакы раз преходити к колцу дробов по \(p\), при чем \(p\) бежи чрез все просте идеалы из \(\Omega\). Ако \(\mathcal C\) јест произвольны \(\oo\)-модул из \(K\), тогда како компонента \(\mathcal C_p\) означаје се разширјены модул \(\mathcal C\oo_p\). Тогда држи:

\paragraph{Лема.}
\(\mathcal C\) јест пресек всех разширјеных модулов \(\mathcal C_p\). Бо ако \(w\) лежи в \(\mathcal C_p\), тогда јест \(\sigma\in\oo\), не делимо чрез \(p\), тако же \(\sigma w\in\mathcal C\). Ако \(w\) лежи в пресеку всех \(\mathcal C_p\), и \(g\) означаје идеал всех \(\sigma\in\oo\) с \(\sigma w\in\mathcal C\), тогда \(g\) не може быти делимы чрез никако \(p\); теды \(g=\oo\). Лема не предпоклада највышши ранг модула.

\paragraph{Последице.}
Всакы модул јест определьен својими компонентами. Ако \(\mathcal D\) јест прави надмодул од \(\mathcal C\), тогда најменеј једно \(\mathcal D_p\) јест прави надмодул од \(\mathcal C_p\). Особливо, ако \(\mathcal M\) јест порјадок в \(K\) и все компоненты \(\mathcal M_p\) сут максималне порјадкы по \(\oo_p\), тогда и \(\mathcal M\) јест максималны порјадок.

\paragraph{Доказ теорема 1.}
\(\mathcal O_a\) јест порјадок, бо \(\mathcal O_a\) јест конечны \(\oo\)-модул највышшего ранга. Далеј
\[
        \mathcal O_a^2
        =\bar aE_\Gg a\,\bar aE_\Gg a
        =\bar aE_\Gg a\,\Sp(a\bar a)
        \subseteq \bar aE_\Gg a=\mathcal O_a.
\]
Да показати максималност, по последици достаточно јест показати максималност всакој компоненты. Како \(\oo_p\) јест колцо главных идеалов, \(a_p\) и \(\bar a_p\) имајут независне комплементарне базы \(a_1,\ldots,a_n\) и \(\bar a_1,\ldots,\bar a_n\). Продукты \(\bar a_iE_\Gg a_k\) творе по §1 матричне јединице; теды
\[
        \mathcal O_{a,p}=\sum_{i,k}(\bar a_iE_\Gg a_k)\oo_p
\]
јест максималны. Како \(\mathcal O_a\) јест максималны, следује такодје \(\Sp(a\bar a)=\oo\).

\paragraph{Доказ теорема 2.}
Из правила следа следује:
\[
\begin{aligned}
 \mathcal O\mathcal L&=oE_\Gg o\,oE_\Gg a=oE_\Gg a=\mathcal L,\\
 \mathcal R\mathcal O&=\bar aE_\Gg o\,oE_\Gg o=\bar aE_\Gg o=\mathcal R,\\
 \mathcal L\mathcal R&=oE_\Gg a\,\bar aE_\Gg o=oE_\Gg o=\mathcal O,\\
 \mathcal R\mathcal L&=\bar aE_\Gg o\,oE_\Gg a=\bar aE_\Gg a=\mathcal O_a.
\end{aligned}
\]

\paragraph{Доказ теорема 3.}
Нехај \(\mathcal L\) јест левы идеал по \(\mathcal O\). Тогда особливо
\[
        \mathcal L=oE_\Gg\mathcal L
        =(oE_\Gg l_1,\ldots,oE_\Gg l_n)
\]
с \(l_1,\ldots,l_n\) како \(\oo\)-базоју од \(\mathcal L\). Положивши \(l_\mu=\sum_i u_{S_i}z_{i\mu}\), добивамо \(E_\Gg l_\mu=E_\Gg a_\mu\), теды \(\mathcal L=oE_\Gg a\). При регуларном \(\mathcal L\) мора \(a\) имети највышши ранг \(n\); тогда \(\bar aE_\Gg o\) јест реципрочны правы идеал. Обча тврдньа следује чрез композицију.

\paragraph{Доказ теорема 4.}
Пресек \(t=[k,\mathcal O_a]\) јест порјадок. Како подмножство од \(\mathcal O_a\), он ставаје се областју правых мултипликаторов и јест највечши порјадок из \(o\), кторы тој услов изполньа; зато сигурно всебује порјадок \(o_a\) од \(a\). Да показати \(t=o_a\), достаточна јест компонентна равност. Ако \(\tau\in t_p\), тогда \(\tau c_{ik}\) за всако \(c_{ik}=\bar a_iE_\Gg a_k\) јест линеарна комбинација \(c_{ij}\) с коефициентами из \(\oo_p\). S другеј страны \(a_k\tau\) ставаје се линеарноју комбинацијеју \(a_j\) с коефициентами из \(\oo_p\). Заради јединства представјеньа чрез \(c_{ij}\) те коефициенты морајут лежати в \(\oo_p\); теды \(\tau\) лежи в порјадку од \(a_p\). Тым \(t\) јест равны порјадку од \(a\).

\paragraph{Доказ теорема 4а.}
Теорем јест специализација теорема 4 на главну област. Ако најпрво \(\mathcal O=oE_\Gg o\) и \(o\) јест такодје в правом порјадку од \(\mathcal L=oE_\Gg a\), тогда \(\bar aE_\Gg a\) належи к главној области; теды \(a\) јест идеал из \(k\), и \(\mathcal L=\mathcal O a\). Ако изходимо из \(\mathcal O_a\), одповедајуче \(\mathcal L=\bar aE_\Gg b\), и \(a^{-1}b\) јест идеал из \(k\); добивамо \(\mathcal L=\mathcal O_a a^{-1}b\).

\subsection*{§III. Максималне порјадкы произвольных скрежених продуктов}

Теоремы изведене за разложене скрежене продукты остајут правдиве и в обчем случају, когда се разделба на области уточньује. Факты о главној области тогда држе точно, друге морајут быти формуловане како тврдјеньа о појединых местах. При том место ныне треба разумети како \(p\)-адично разширенье.

\paragraph{Дефиниција.}
Нехај \(K\) јест проста нормална алгебра над \(\Omega\), а \(k\) максимално комутативно подполье. В област по \(k\) рачут се все максималне порјадкы из \(K\), кторе имајут једнакы пресек с \(k\) и кроме того на разветвјених местах од \(K\) имајут једнаке \(p\)-адичне компоненты. Ако пресек с \(k\) јест главны порјадок, тогда иде о главну област.

За главне области тогда држи:

\paragraph{Теорем 1.}
Максималне порјадкы једнеј главној области преходет једне в друге трансформацијеју с разширеными идеалами идеалов из \(k\). Другыми словами: леве идеалы максималных порјадков главној области, кторе сут приме к диференте и всебујут \(o\) в својем правом порјадку, сут разширеньа идеалов из \(k\).

\paragraph{Теорем 2.}
Ако \(K\) јест простого степена, тогда јест само једна главна област; она преходи в себе чрез разширјене идеалы \(k\)-идеалов. Одповедајучи идеалы фиксного максималного порјадка \(\mathcal O\) сут идеалы из \(k\), разширјене з двустранными идеалами из \(\mathcal O\).

\paragraph{Доказ.}
Два максималне порјадкы једнеј главној области различе се само на конечно много местах; по дефиницији то сут места, на кторых \(K\) разклада се. Тут \(K\) има форму разгледану в §§1--2. То јест јасно, ако \(k\) јест Галоисово; ако факторна система јест асоциована с јединицоју, може быти заменьена јединичноју системоју. По §1 то држи и в не-Галоисовом случају. Теды на всаком таком месту и зато обче иде о разширјене идеалы.

Ако \(K\) јест простого степена, она јест или разложена или делитевна алгебра. В другом случају она остаје делитевна алгебра на всех разветвјених местах и зато има само једну главну област, также знову држи трансформација с разширеными идеалами. Најобчејша трансформација добива се компонованьем разширјеных идеалов с идеалами, кторе трансформујут максималны порјадок в себе, теды с двустранными идеалами; они се, како јест знано, разликујут од централных идеалов само на разветвјених местах, теды од разширјеных идеалов.

Далеј држи:

\paragraph{Теорем 3.}
Максималне порјадкы произвольној области преходет једне в друге чрез идеалы приме к диференте, кторе на всаком месту сут составјене из модулов одповедајучего порјадка тако како указано в §2, теорем 3. Компоненты тых модулов при фиксном идеалу различе се од порјадка само на конечно много местах.

То следује непосредно из §2; треба само имети на памети, же операторы \(u_S\) на појединых местах морајут быти заменьене одповедајучими еквивалентными.

Чи и тут к всакому порјадку из \(k\) сут максималне порјадкы с тым порјадком како пресеком, потребује ближше изследованье на разветвјених местах. Главна област всегда јест, како можно взяти из скреженого, евентуално обобченого, продуктного представјеньа, кторе води к порјадку всебујучему \(o\), когда факторне системы сут взяте како целе.

Далејши вопрос јест, в кој мери можно аритметично разделбо по областах карактеризовати алгебраичне разпадне польа. Другыми словами: производе различне разпадне польа всегда различне областне разделбы, или уж различне главне области? Изчерпају ли можно главне области всех разпадных поль всу совокупност максималных порјадков?

Же така карактеризација чрез једен самы максималны порјадок сигурно неможна јест, показују најпростејши примеры. Тако максималны порјадок
\[
        \left[\,1,i,j,\frac{1+i+j+k}{2}\,\right]
\]
кватернионского тела всебује поред главного порјадка \([1,i]\) такодје главны порјадок
\[
        \left[\,1,\frac{1+i+j+k}{2}\,\right]
\]
польа третјих кореньев јединице.

\begin{center}
Готтинген, август 1932.
\end{center}
\clearpage

\end{document}

